\documentclass[10pt,a4paper]{amsart}
\usepackage{amsmath,amssymb,amsthm,graphicx}
% Retained preamble for the Stacks Project chapter snapshot in this directory.
% The chapter's own preamble.tex is a copy of the shared Stacks Project preamble; the block
% below is sliced MECHANICALLY, line for line, from it. Nothing was added, removed or
% redefined; the wrapper class, the xr-hyper/hyperref wiring and the externaldocument list
% are omitted because the assembler supplies the document class.
%
% lines taken from preamble.tex: 3, 5, 9, 11, 12, 156, 157, 158, ... (51 lines total)
%
% package coverage restored from the chapter's own preamble (2026-10-13): xy, xy-2cell

\usepackage[all]{xy}
\xyoption{2cell}
\UseAllTwocells

\newenvironment{reference}{\comment}{\endcomment}
\newenvironment{slogan}{\comment}{\endcomment}
\newenvironment{history}{\comment}{\endcomment}
\theoremstyle{plain}
\newtheorem{theorem}[subsection]{Theorem}
\newtheorem{proposition}[subsection]{Proposition}
\newtheorem{lemma}[subsection]{Lemma}
\theoremstyle{definition}
\newtheorem{definition}[subsection]{Definition}
\newtheorem{example}[subsection]{Example}
\newtheorem{exercise}[subsection]{Exercise}
\newtheorem{situation}[subsection]{Situation}
\theoremstyle{remark}
\newtheorem{remark}[subsection]{Remark}
\newtheorem{remarks}[subsection]{Remarks}
\numberwithin{equation}{subsection}
\def\lim{\mathop{\mathrm{lim}}\nolimits}
\def\colim{\mathop{\mathrm{colim}}\nolimits}
\def\Spec{\mathop{\mathrm{Spec}}}
\def\Hom{\mathop{\mathrm{Hom}}\nolimits}
\def\Ext{\mathop{\mathrm{Ext}}\nolimits}
\def\SheafHom{\mathop{\mathcal{H}\!\mathit{om}}\nolimits}
\def\SheafExt{\mathop{\mathcal{E}\!\mathit{xt}}\nolimits}
\def\Sch{\mathit{Sch}}
\def\Mor{\mathop{\mathrm{Mor}}\nolimits}
\def\Ob{\mathop{\mathrm{Ob}}\nolimits}
\def\Sh{\mathop{\mathit{Sh}}\nolimits}
\def\NL{\mathop{N\!L}\nolimits}
\def\CH{\mathop{\mathrm{CH}}\nolimits}
\def\proetale{{pro\text{-}\acute{e}tale}}
\def\etale{{\acute{e}tale}}
\def\QCoh{\mathit{QCoh}}
\def\Ker{\mathop{\mathrm{Ker}}}
\def\Im{\mathop{\mathrm{Im}}}
\def\Coker{\mathop{\mathrm{Coker}}}
\def\Coim{\mathop{\mathrm{Coim}}}
\DeclareMathSymbol{\boxtimes}{\mathbin}{AMSa}{"02}
\def\QCohstack{\mathcal{QC}\!\mathit{oh}}
\def\Cohstack{\mathcal{C}\!\mathit{oh}}
\def\Spacesstack{\mathcal{S}\!\mathit{paces}}
\def\Quotfunctor{\mathrm{Quot}}
\def\Hilbfunctor{\mathrm{Hilb}}
\def\Curvesstack{\mathcal{C}\!\mathit{urves}}
\def\Polarizedstack{\mathcal{P}\!\mathit{olarized}}
\def\Complexesstack{\mathcal{C}\!\mathit{omplexes}}
\def\Pic{\mathop{\mathrm{Pic}}\nolimits}
\def\Picardstack{\mathcal{P}\!\mathit{ic}}
\def\Picardfunctor{\mathrm{Pic}}
\def\Deformationcategory{\mathcal{D}\!\mathit{ef}}
\usepackage{cleveref}
\crefname{originallemma}{original lemma}{original lemmas}
\Crefname{originallemma}{Original lemma}{Original lemmas}
\makeatletter
\newlabel{lemma-functoriality-sheaves}{{lemma-functoriality-sheaves}{source}{}{sourcepointer-1}{}}
\newlabel{lemma-functoriality-sheaves@cref}{{[originallemma][0][]lemma-functoriality-sheaves}{[1][0][]source}{}{}{}}
\makeatother
\crefname{originallemma}{original lemma}{original lemmas}
\Crefname{originallemma}{Original lemma}{Original lemmas}
\makeatletter
\newlabel{sites-section-continuous-functors}{{00WU}{source}{}{sourcepointer-2}{}}
\newlabel{sites-section-continuous-functors@cref}{{[originallemma][0][]00WU}{[1][0][]source}{}{}{}}
\makeatother
\crefname{originallemma}{original lemma}{original lemmas}
\Crefname{originallemma}{Original lemma}{Original lemmas}
\makeatletter
\newlabel{sites-section-functoriality-PSh}{{00VC}{source}{}{sourcepointer-3}{}}
\newlabel{sites-section-functoriality-PSh@cref}{{[originallemma][0][]00VC}{[1][0][]source}{}{}{}}
\makeatother
\title{Stacks, Sheaves on Algebraic Stacks: lemma-representable-pushforward}
\author{The Stacks Project Authors}
\date{stacks-project:stacks-sheaves.tex:local-snapshot-20261012; GNU FDL 1.2 or later}
\begin{document}
\maketitle
\noindent\textit{Editorial note.} Statement span is lines 873-886 of sources/stacks-algebraic-stacks-sheaves/pinned-original-source.tex; proof span is lines 888-914, the complete author proof (25 body lines).\par\medskip
\noindent\textit{Editorial note.} Navigation stubs: 3 cited-but-not-included labels are supplied as non-load-bearing pointers with their pinned tags where the snapshot tag index has them.\par\medskip
\noindent\textit{Editorial note.} extra\_packages is empty on purpose: the retained preamble already loads xy with options and a second declaration clashes in the assembler wrapper.\par\medskip
\noindent\textit{Editorial note.} Text is retained verbatim, including the author's own typos and placeholders.\par\medskip
\noindent\textit{Editorial note.} Queue index 23 of queue\_algebraic-stacks-sheaves.json (pb=1162, stmt 873-886, proof 888-914).\par\medskip
\noindent\textit{Allowed standard prerequisites (editorial).}\begin{itemize}
\item Explicit prerequisite boundary. The cited inputs are prerequisites; the argument in the retained span is not replaced by any cited result.
\end{itemize}\medskip
\ifdefined\hypertarget\hypertarget{sourcepointer-1}{}\fi
\noindent\textit{Source reference.} Original lemma lemma-functoriality-sheaves: sources/stacks-algebraic-stacks-sheaves/pinned-original-source.tex, source lines 678--686. Cited by the retained statement or proof but not part of the collected spans; supplied only as a navigation pointer. The argument does not depend on its content.\par\medskip
\ifdefined\hypertarget\hypertarget{sourcepointer-2}{}\fi
\noindent\textit{Source reference.} Original lemma 00WU: sources/stacks-algebraic-stacks-sheaves/pinned-original-source.tex, source lines 914--914. Cited by the retained statement or proof but not part of the collected spans; supplied only as a navigation pointer. The argument does not depend on its content.\par\medskip
\ifdefined\hypertarget\hypertarget{sourcepointer-3}{}\fi
\noindent\textit{Source reference.} Original lemma 00VC: sources/stacks-algebraic-stacks-sheaves/pinned-original-source.tex, source lines 914--914. Cited by the retained statement or proof but not part of the collected spans; supplied only as a navigation pointer. The argument does not depend on its content.\par\medskip
\begin{lemma}
\label{lemma-representable-pushforward}
Let $f : \mathcal{X} \to \mathcal{Y}$ be a representable $1$-morphism of
categories fibred in groupoids over $(\Sch/S)_{fppf}$. Let
$\tau \in \{Zar, \etale, smooth, syntomic, fppf\}$.
Then the functor $u : \mathcal{Y}_\tau \to \mathcal{X}_\tau$ is continuous
and defines a morphism of sites $\mathcal{X}_\tau \to \mathcal{Y}_\tau$
which induces the same morphism of topoi
$\Sh(\mathcal{X}_\tau) \to \Sh(\mathcal{Y}_\tau)$
as the morphism $f$ constructed in
Lemma \ref{lemma-functoriality-sheaves}.
Moreover, $f_*\mathcal{F}(y) = \mathcal{F}(u(y))$ for any presheaf
$\mathcal{F}$ on $\mathcal{X}$.
\end{lemma}

\begin{proof}
Let $\{y_i \to y\}$ be a $\tau$-covering in $\mathcal{Y}$. By definition
this simply means that $\{q(y_i) \to q(y)\}$ is a $\tau$-covering of
schemes. By the final remark above the lemma we see that
$\{p(u(y_i)) \to p(u(y))\}$ is the base change of the $\tau$-covering
$\{q(y_i) \to q(y)\}$ by $p(u(y)) \to q(y)$, hence is itself a
$\tau$-covering by the axioms of a site. Hence $\{u(y_i) \to u(y)\}$
is a $\tau$-covering of $\mathcal{X}$. This proves that $u$ is
continuous.

\medskip\noindent
Let's use the notation $u_p, u_s, u^p, u^s$ of
Sites, Sections \ref{sites-section-functoriality-PSh} and
\ref{sites-section-continuous-functors}.
If we can show the final assertion of the lemma, then we see that
$f_* = u^p = u^s$ (by continuity of $u$ seen above) and hence by adjointness
$f^{-1} = u_s$ which will prove $u_s$ is exact, hence that $u$ determines
a morphism of sites, and the equality will be clear as well.
To see that $f_*\mathcal{F}(y) = \mathcal{F}(u(y))$ note that by
definition
$$
f_*\mathcal{F}(y) = ({}_pf\mathcal{F})(y) =
\lim_{f(x) \to y} \mathcal{F}(x).
$$
Since $u(y)$ is a final object in the category the limit is taken
over we conclude.
\end{proof}

\end{document}
